
\chapter[Implementation]{Implementation}{Implementation}\label{im}

\noindent
\rule{0.49\textwidth}{0.75pt} $_{\bigcirc}$ \rule{0.49\textwidth}{0.75pt}\\


\section{Implementation Tools, Data Handling Approach and System Limitations}
The implementation of DAVS focuses on building an end-to-end secure federated learning pipeline for medical imaging. The training workflow is implemented as a coordinator-client architecture in which each hospital client trains locally and communicates only model updates. To reduce bandwidth, each client applies Random Projection-based gradient sketching to compress high-dimensional updates into low-dimensional sketches before transmission.

Data handling follows federated principles: raw medical images remain within each institution, and only compressed update representations and evaluation metrics are exchanged. The coordinator aggregates client updates (FedAvg baseline) and then executes DAVS to select a validator committee using current-round sketch statistics (cosine similarity and magnitude consistency). The selected committee runs a PBFT-inspired validation procedure to accept or reject the proposed update. Once validated, a permissioned blockchain audit layer records round metadata such as model hashes, committee IDs, DAVS scores, PBFT votes, and performance metrics.

Limitations:
The implementation is constrained by the scale and realism of the experimental setup. Communication and computation trade-offs depend on sketch dimension choices and the number of clients. While gradient sketching provides strong bandwidth reduction, it introduces an approximation layer that must be calibrated to preserve similarity reliably. PBFT-style validation adds overhead proportional to committee size and message exchange. The permissioned blockchain layer provides immutable logging, but long-term storage growth and integration with real clinical governance processes remain future work. Additionally, the experimental evaluation reflects common medical FL settings (non-IID partitions and synthetic attack models) rather than full production deployments.

\section{Identity Data Structure and Cryptographic Components Description}
Core Training and Audit Data:

Client Training Data:
Local medical imaging datasets held privately by each hospital client (non-IID across institutions).
Applications: Used only for local training; never transmitted to the coordinator or validators.

Model Update Data:
Per-client model updates \(\Delta w_i\) produced after local epochs.
Applications: Input to aggregation and validation; compressed for communication efficiency.

Gradient Sketch Data:
Low-dimensional sketches \(s_i = R \cdot \Delta w_i\) generated using Random Projection.
Applications: Used for DAVS scoring (cosine similarity and magnitude consistency) and lightweight trust evaluation.

Consensus and Audit Data:
PBFT vote logs, committee identifiers, global model hash per round, DAVS scores, and evaluation metrics.
Applications: Stored in the permissioned blockchain ledger to provide immutable training provenance and enable forensic audit.



\section{Cryptographic Data Preparation and Processing}

DAVS substitutes heavy cryptographic preprocessing with lightweight, privacy-preserving update processing tailored for federated learning. After local training, each client computes a high-dimensional update and then applies Random Projection to produce a compact sketch. This sketching step reduces communication cost and obfuscates the raw gradient structure while preserving similarity properties required for DAVS scoring.

For provenance, the coordinator computes hashes of the aggregated global model per round. These hashes, together with committee membership, scores, votes, and metrics, are written to the permissioned ledger. This provides integrity and tamper-evidence for the training process without storing raw data or requiring trusted external parties.

\section{Description of Algorithms (Modular)}

Federated Averaging (FedAvg):\\
Algorithm: Standard FL aggregation by averaging client updates.\\
Usage: Establishes baseline federated training and provides the initial aggregation step before validation.

Random Projection Gradient Sketching:\\
Algorithm: Random Projection Compression (RPC).\\
Usage: Compresses high-dimensional updates into low-dimensional sketches that preserve cosine similarity and reduce communication overhead.

Data-Aware Verifier Selection (DAVS):\\
Algorithm: Amnesic scoring based on cosine similarity and magnitude consistency on sketch vectors.\\
Usage: Selects the top-\(k\) representative clients as a validator committee and filters out anomalous/malicious updates.

PBFT-inspired Consensus Validation:\\
Algorithm: Practical Byzantine Fault Tolerance style voting (Pre-Prepare, Prepare, Commit).\\
Usage: Validators check update quality and vote to accept/reject the proposed global update; global model updates only on quorum.

Permissioned Blockchain Audit Logging:\\
Algorithm: Hash-chained ledger with per-round metadata.\\
Usage: Records model hashes, committee membership, DAVS scores, PBFT vote logs, and metrics for tamper-proof provenance.

\section{Description of Tools for Implementation}

VS Code Visual Studio code:
Description: A source code editor created by Microsoft, which is light in weight but powerful.
Usage: Writing, debugging and administration of the federated learning pipeline, validation logic, and audit ledger code.

Git:
Description: A distributed version control system to track source code changes.
Usage: Controls versioning of an entire system during development, facilitating effective collaboration, tracking of code and retracting changes during various phases of the project.

GitHub:
Description: A web based platform for hosting and collaborating on Git repositories.
Usage: Manages the DAVS project repository where people can collaborate, manage versions, track issues and maintain the project history with adequate documentation.

Postman:
Description: An API development and testing tool.
Usage: (Optional) Tests and debugs any coordinator/client API endpoints if exposed as services.

MetaMask:
Description: A blockchain interaction and cryptocurrency wallet. Applications: Authenticate wallets, sign transactions and allow users to safely communicate with Ethereum smart contracts when verifying their identity and controlling access in an identity verification and access control operation.
Description: Not required for the DAVS training pipeline; the audit layer operates in a permissioned environment without user-facing wallet interaction.


\section{Interface Design}
The interface for DAVS is primarily experiment- and monitoring-oriented. A lightweight monitoring view (or console logs) can be used to track round progression, client participation, sketching compression statistics, committee selection outcomes, PBFT votes, model acceptance/rejection, and accuracy/loss curves. The focus is on clarity of training provenance and verification outcomes rather than user registration workflows.


\rule{0.49\textwidth}{0.75pt} $_{\bigcirc}$ \rule{0.49\textwidth}{0.75pt}\\
